% !TEX program = tectonic
% Standalone LaTeX edition. Content and section order follow the matching Markdown.
% Edit this file for typesetting; keep any later content revision synchronized with Markdown.
\documentclass[UTF8,fontset=none,zihao=-4,a4paper]{ctexart}
\usepackage{waterhand-report}
\renewcommand{\ReportTitle}{实验补充说明}
\renewcommand{\thetable}{S\arabic{table}}
\hypersetup{pdftitle={WaterHand Processor Development Skills 实验补充说明},pdfauthor={计算机科学与技术学院 付家乐},pdfsubject={产品版本 v3.0.3}}
\begin{document}
\setstretch{1.15}
\small
\ctexset{
  section={beforeskip=13pt,afterskip=7pt},
  subsection={beforeskip=10pt,afterskip=5pt}
}
\setlength{\LTpre}{5pt}
\setlength{\LTpost}{5pt}
\begin{center}
{\sffamily\bfseries\large WaterHand Processor Development Skills\par}
\vspace{5mm}
{\heiti\LARGE 实验补充说明\par}
\vspace{3mm}
{\small v3.0.3\quad 2026 年 9 月 6 日\par}
\end{center}
\vspace{3mm}
本文随\href{DESIGN_REPORT.pdf}{《设计文档》}提交，集中说明三项实验的运行条件、完整指标及验证范围。设计文档第 5 章给出评估目的、主要结果与有效性判断；本说明提供理解和复核这些结果所需的细节。

\section{实验条件与统计口径}

\subsection{角色设置与组别隔离}

顺序双发射开发对照和分支预测增量开发对照均使用同等配置的 Codex 主会话担任代理设计师，在给定目标内选择方案、组织工程子任务并复核产出。这一安排用于控制人类工程师之间的经验差异和重复任务中的学习效应，实验结果同时包含方案选择与工程执行的影响。第三项实验由真实工程师主导，展示产品预期的协作方式。

两项对照均按 Skill 后 Control 的顺序串行运行。两组使用相同的公共任务、项目起点、工具条件和通用 Skill，分别维护 Git 仓库、Codex 配置与会话目录、Memory 和结果目录。Memory 的读取与生成均启用，从相同的空白内容开始，各组副本独立可写。

Skill 组的主会话与工程子任务可使用冻结的产品 Skill Package，Control 组的对应角色均禁止读取本产品，包括全局安装的副本。Skill 带来的额外上下文计入该组资源消耗。

\subsection{模型与工具配置}

两项对照的主任务均使用 \code{gpt-\allowbreak{}5.\allowbreak{}6-\allowbreak{}sol}、\code{ultra} 推理设置，工程子任务使用同一模型的 \code{high} 设置，同时运行的子任务上限为 4，累计调用可超过该值。实验侧未另设 Token、实际经过时间或主轮次上限。人工接管后的主会话使用 \code{gpt-\allowbreak{}5.\allowbreak{}6-\allowbreak{}sol}、\code{high}。

分支预测对照使用 Codex 0.151.0 和冻结的 Windows 工具链，Skill 组获得六项产品 Skill，组织者未注入技术设计决定。两组继承顺序双发射第二次对照中同一份 Skill 实现及其文档、工具脚本，输入的预测草稿相同，未获得人工案例后续确定的设计答案。

\subsection{指标含义}

性能使用 CoreMark 的 performance 和 validation 两个工作负载。IPC 为退休指令数除以测量周期数 \code{measuredCycles}。本文的性能比较限于 RTL 仿真，物理频率、面积与功耗均未测量。

Token 为模型请求计数，缓存输入已包含在输入中，推理输出已包含在输出中，均不重复累加。主会话计数与可识别工程主、子会话累计分别注明，后者不含单列的组织者和自动审批资源。Token 数量不直接等于账单金额或人类认知成本。

任务用时反映各次运行的自行交付过程，具体扣除口径在相应实验中说明。任务用时与最终验收状态分别统计，并结合交付结果比较。

\section{顺序双发射处理器开发对照}

\subsection{验收与仿真性能}

两次对照从相同的架构目标和构建骨架出发，要求保持顺序双发射硬目标、受保护的七级流水角色、RV32I 对外行为、固定顶层接口 \code{ExperimentTop} 及两个 CoreMark 镜像。组织者对四个交付实现的独立副本执行相同的 13 个定向场景和两个 CoreMark 工作负载，四个实现均通过全部 15 项检查。

表~\ref{tab:S1} 给出各工作负载的 IPC。两个保留样本中，Control 均以粗粒度串行化减少跨组关联状态，Skill 均通过跨级数据转发支持不同指令组在流水线中重叠执行。相同工作负载的退休指令数一致，IPC 差异来自执行周期数。

\begingroup
\fontsize{9.5}{14}\selectfont
\setstretch{1}
\renewcommand{\arraystretch}{1.15}
\setlength{\ReportTableWidth}{\dimexpr\linewidth-8\tabcolsep\relax}
\begin{longtable}{@{}>{\raggedright\arraybackslash}p{0.17\ReportTableWidth}>{\centering\arraybackslash}p{0.21\ReportTableWidth}>{\centering\arraybackslash}p{0.18\ReportTableWidth}>{\centering\arraybackslash}p{0.19\ReportTableWidth}>{\centering\arraybackslash}p{0.25\ReportTableWidth}@{}}
\caption{顺序双发射对照的 CoreMark IPC}\label{tab:S1}\\
\toprule
{\bfseries 运行} & {\bfseries Workload} & {\bfseries Skill IPC} & {\bfseries Control IPC} & {\bfseries Skill 相对增幅} \\
\midrule
\endfirsthead
\multicolumn{5}{c}{\small 表~\thetable\quad 顺序双发射对照的 CoreMark IPC（续）}\\[4pt]
\toprule
{\bfseries 运行} & {\bfseries Workload} & {\bfseries Skill IPC} & {\bfseries Control IPC} & {\bfseries Skill 相对增幅} \\
\midrule
\endhead
\midrule
\multicolumn{5}{r}{\footnotesize 接下页}\\
\endfoot
\bottomrule
\endlastfoot
第一次对照 & performance & 0.285884 & 0.211433 & 35.21\% \\[3pt]
第一次对照 & validation & 0.284909 & 0.211149 & 34.93\% \\[3pt]
第二次对照 & performance & 0.303263 & 0.206840 & 46.62\% \\[3pt]
第二次对照 & validation & 0.303437 & 0.206836 & 46.70\% \\[3pt]
\end{longtable}
\endgroup

\subsection{开发用时与资源消耗}

表~\ref{tab:S2} 汇总首次 CoreMark 通过时间、完整任务用时及主会话 Token 消耗。时间扣除仅覆盖日志可直接归因的 sbt 锁区间；并行工作中这些区间未必全部位于完成关键路径，因此校正值属于反事实估计。

\begingroup
\fontsize{9.5}{14}\selectfont
\setstretch{1}
\renewcommand{\arraystretch}{1.15}
\setlength{\ReportTableWidth}{\dimexpr\linewidth-4\tabcolsep\relax}
\begin{longtable}{@{}>{\raggedright\arraybackslash}p{0.32\ReportTableWidth}>{\raggedright\arraybackslash}p{0.34\ReportTableWidth}>{\raggedright\arraybackslash}p{0.34\ReportTableWidth}@{}}
\caption{顺序双发射对照的开发用时与主会话 Token 消耗}\label{tab:S2}\\
\toprule
{\bfseries 指标} & {\bfseries 第一次对照 Skill / Control} & {\bfseries 第二次对照 Skill / Control} \\
\midrule
\endfirsthead
\multicolumn{3}{c}{\small 表~\thetable\quad 顺序双发射对照的开发用时与主会话 Token 消耗（续）}\\[4pt]
\toprule
{\bfseries 指标} & {\bfseries 第一次对照 Skill / Control} & {\bfseries 第二次对照 Skill / Control} \\
\midrule
\endhead
\midrule
\multicolumn{3}{r}{\footnotesize 接下页}\\
\endfoot
\bottomrule
\endlastfoot
首次 CoreMark 通过，扣除可归因 sbt 锁区间 & 24 分 58.1 秒 / 21 分 44.8 秒 & 37 分 32.5 秒 / 56 分 26.1 秒 \\[3pt]
完整任务用时，同一扣除口径 & 69 分 40.7 秒 / 71 分 21.7 秒 & 65 分 59.5 秒 / 61 分 2.1 秒 \\[3pt]
主会话累计总 Token & 31,430,282 / 29,589,081 & 32,274,107 / 24,335,568 \\[3pt]
\end{longtable}
\endgroup

两次对照的主会话 Token 消耗分别增加约 6.22\% 和 32.62\%。第二次对照的可识别工程团队累计为 Skill 68,495,964、Control 52,889,664 Token，增加约 29.51\%。Skill 组两次均取得更高 IPC，完整任务用时的优劣方向并不一致。

\subsection{Skill 使用与工程组织}

第二次对照的主任务及工程角色读取 \code{design-\allowbreak{}chisel-\allowbreak{}processor} 和 \code{implement-\allowbreak{}chisel-\allowbreak{}processor}。

文档角色读取 \code{design-\allowbreak{}chisel-\allowbreak{}processor} 和 \code{organize-\allowbreak{}processor-\allowbreak{}docs}，并运行文档检查器。交付材料形成模块、接口和生命周期等阅读入口，源码按 frontend、backend、core 等责任组织。

两组的公共任务均要求模块拆分和源码旁摘要，实现与文档工作可并行开展。本节对文档组织方法的评价采用第二次对照中记录的 \code{organize-\allowbreak{}processor-\allowbreak{}docs} 调用及相应产出。

\section{分支预测增量开发对照}

\subsection{功能与工程一致性}

两组基于相同的可运行处理器和预测草稿，补全预测、恢复、训练及停顿处理。草稿规定 128 项分支历史表（BHT）和 7 位全局分支历史状态（GHT），要求保持既有流水线和双发射能力。功能检查、正式文档一致性和完整验收结果见表~\ref{tab:S3}。

\begingroup
\fontsize{9.5}{14}\selectfont
\setstretch{1}
\renewcommand{\arraystretch}{1.15}
\setlength{\ReportTableWidth}{\dimexpr\linewidth-4\tabcolsep\relax}
\begin{longtable}{@{}>{\raggedright\arraybackslash}p{0.40\ReportTableWidth}>{\raggedright\arraybackslash}p{0.28\ReportTableWidth}>{\raggedright\arraybackslash}p{0.32\ReportTableWidth}@{}}
\caption{分支预测对照的功能检查与完整验收结果}\label{tab:S3}\\
\toprule
{\bfseries 指标} & {\bfseries Skill} & {\bfseries Control} \\
\midrule
\endfirsthead
\multicolumn{3}{c}{\small 表~\thetable\quad 分支预测对照的功能检查与完整验收结果（续）}\\[4pt]
\toprule
{\bfseries 指标} & {\bfseries Skill} & {\bfseries Control} \\
\midrule
\endhead
\midrule
\multicolumn{3}{r}{\footnotesize 接下页}\\
\endfoot
\bottomrule
\endlastfoot
独立外部场景 & 24/24 通过 & 24/24 通过 \\[3pt]
团队新增测试独立复跑 & 8/8 通过 & 7/7 通过 \\[3pt]
组织者追加探针 & 2/2 通过 & 2/2 通过 \\[3pt]
BP01 至 BP08 & 通过 & 通过 \\[3pt]
BP09 文档与实现一致性 & 通过 & 失败 \\[3pt]
最终完整验收 & 通过 & 未通过，文档一致性缺陷 \\[3pt]
\end{longtable}
\endgroup

BP01 至 BP09 是任务的九项验收要求。最终匿名审查确认，Control 的 Architecture 仍规定已经不存在的 \code{redirect\_\allowbreak{}applied} 状态和集中事件仲裁，与新增前端本地预测重定向及实际恢复行为冲突。

24 项外部场景检查程序行为，新增预测功能由源码审查、预测状态测试和组织者探针共同核对。预测状态探针覆盖 4096 周期、3141 次训练更新和 877 次共同主查询同址碰撞；整核探针使用四个随机种子检查背压、延迟响应、运行中复位、退休结果及错误路径副作用。

Control 还提供了整核预测、解析和训练诊断，并覆盖寄存器间接跳转 JALR 的目标等于顺序地址等边界。两组对较老异常、待采用预测和旧响应同拍竞争的组合场景均未进行直接测试。

\subsection{仿真性能}

表~\ref{tab:S4} 给出共同基线与两组实现的 CoreMark IPC。两组相对基线均提高约 13.5\% 至 13.9\%，彼此差距很小。

\begingroup
\fontsize{9.5}{14}\selectfont
\setstretch{1}
\renewcommand{\arraystretch}{1.15}
\setlength{\ReportTableWidth}{\dimexpr\linewidth-6\tabcolsep\relax}
\begin{longtable}{@{}>{\raggedright\arraybackslash}p{0.22\ReportTableWidth}>{\centering\arraybackslash}p{0.26\ReportTableWidth}>{\centering\arraybackslash}p{0.26\ReportTableWidth}>{\centering\arraybackslash}p{0.26\ReportTableWidth}@{}}
\caption{分支预测对照的 CoreMark IPC}\label{tab:S4}\\
\toprule
{\bfseries Workload} & {\bfseries 共同基线 IPC} & {\bfseries Skill IPC} & {\bfseries Control IPC} \\
\midrule
\endfirsthead
\multicolumn{4}{c}{\small 表~\thetable\quad 分支预测对照的 CoreMark IPC（续）}\\[4pt]
\toprule
{\bfseries Workload} & {\bfseries 共同基线 IPC} & {\bfseries Skill IPC} & {\bfseries Control IPC} \\
\midrule
\endhead
\midrule
\multicolumn{4}{r}{\footnotesize 接下页}\\
\endfoot
\bottomrule
\endlastfoot
performance & 0.303263 & 0.344558 & 0.344247 \\[3pt]
validation & 0.303437 & 0.345687 & 0.345561 \\[3pt]
\end{longtable}
\endgroup

Skill 两个工作负载分别比 Control 少 2041 和 879 个测量周期，周期差约为 0.0901\% 和 0.0364\%。三个带背压的公开定向程序中，两组均比基线多用 7 拍。

\subsection{开发用时与资源消耗}

表~\ref{tab:S5} 记录两组自行交付用时及可识别工程会话的 Token 消耗。本项采用原始工程用时，未扣除基础设施等待。

\begingroup
\fontsize{9.5}{14}\selectfont
\setstretch{1}
\renewcommand{\arraystretch}{1.15}
\setlength{\ReportTableWidth}{\dimexpr\linewidth-4\tabcolsep\relax}
\begin{longtable}{@{}>{\raggedright\arraybackslash}p{0.40\ReportTableWidth}>{\raggedright\arraybackslash}p{0.30\ReportTableWidth}>{\raggedright\arraybackslash}p{0.30\ReportTableWidth}@{}}
\caption{分支预测对照的开发用时与工程会话 Token 消耗}\label{tab:S5}\\
\toprule
{\bfseries 指标} & {\bfseries Skill} & {\bfseries Control} \\
\midrule
\endfirsthead
\multicolumn{3}{c}{\small 表~\thetable\quad 分支预测对照的开发用时与工程会话 Token 消耗（续）}\\[4pt]
\toprule
{\bfseries 指标} & {\bfseries Skill} & {\bfseries Control} \\
\midrule
\endhead
\midrule
\multicolumn{3}{r}{\footnotesize 接下页}\\
\endfoot
\bottomrule
\endlastfoot
主任务自行交付用时 & 55 分 24.712 秒 & 61 分 40.455 秒 \\[3pt]
可识别工程主、子会话累计 Token & 61,122,251 & 49,079,067 \\[3pt]
\end{longtable}
\endgroup

Skill 组 Token 消耗增加约 24.54\%。Control 在开发过程中更早取得首次候选 CoreMark 通过，退出时仍未满足文档一致性要求。表~\ref{tab:S5} 中的 Control 用时未包含后续修正文档、达到完整验收所需的时间。

本次对照记录到的产品 Skill 调用为 \code{design-\allowbreak{}chisel-\allowbreak{}processor} 和 \code{implement-\allowbreak{}chisel-\allowbreak{}processor}。文档一致性按实验组的整体工程产出评价；\code{organize-\allowbreak{}processor-\allowbreak{}docs} 的使用效果见第 2.3 节及第 4 节人工案例。

\clearpage
\section{人类工程师主导的分支预测开发}

\subsection{架构决定与工程约束}

工程师接管已有处理器后，要求先整理模块文档，再审查其分支预测草稿。Agent 整理未决条件及候选影响，工程师据此选择方案并确认正式设计，随后授权实现。表~\ref{tab:S6} 对照工程师的主要决定及其需要落实的工程条件。

\begingroup
\fontsize{9.5}{14}\selectfont
\setstretch{1}
\renewcommand{\arraystretch}{1.15}
\setlength{\ReportTableWidth}{\dimexpr\linewidth-2\tabcolsep\relax}
\begin{longtable}{@{}>{\raggedright\arraybackslash}p{0.45\ReportTableWidth}>{\raggedright\arraybackslash}p{0.55\ReportTableWidth}@{}}
\caption{人工案例中工程师决定与正式设计约束的对应关系}\label{tab:S6}\\
\toprule
{\bfseries 工程师作出的决定} & {\bfseries 正式设计需要说明的条件} \\
\midrule
\endfirsthead
\multicolumn{2}{c}{\small 表~\thetable\quad 人工案例中工程师决定与正式设计约束的对应关系（续）}\\[4pt]
\toprule
{\bfseries 工程师作出的决定} & {\bfseries 正式设计需要说明的条件} \\
\midrule
\endhead
\midrule
\multicolumn{2}{r}{\footnotesize 接下页}\\
\endfoot
\bottomrule
\endlastfoot
GHT 使用已解析分支历史 & 更新历史的事件，以及异常、取消分支的处理 \\[3pt]
两条并行指令基于同一拍 BHT 和 GHT 状态预测 & 查询时点、结果保存及其与指令组的对应关系 \\[3pt]
在执行级 EX 向访存级 M1 实际推进时训练 & 停顿时保持，实际接受时单次更新 \\[3pt]
异常分支不参与训练 & 异常、取消和指令有效性对训练资格的约束 \\[3pt]
BHT 同步读写，同拍写入不影响查询结果 & 读旧值、结果延迟及前端寄存边界 \\[3pt]
直接跳转 JAL 在译码级 Decode 静态处理 & 动态预测与静态跳转的范围及恢复行为 \\[3pt]
\end{longtable}
\endgroup

正式 \code{Predict.\allowbreak{}md} 记录这些约定，Agent 随后修改预测器与相关流水级，同步维护源码旁摘要并开展验证。针对同表项读写，\code{BranchPredictorSpec.\allowbreak{}scala} 检查同拍查询与更新后，下一拍返回旧值 1；后一次查询调整 PC 以抵消 GHT 变化，再检查同一表项的新值 2。

\subsection{功能检查与仿真结果}

表~\ref{tab:S7} 列出改造前后的周期数、退休指令数和 IPC。前后工作负载镜像校验值和退休指令数一致，保留的五项 Scala 测试及 15 项外部功能场景均通过。

\begingroup
\fontsize{9.5}{14}\selectfont
\setstretch{1}
\renewcommand{\arraystretch}{1.15}
\setlength{\ReportTableWidth}{\dimexpr\linewidth-10\tabcolsep\relax}
\begin{longtable}{@{}>{\raggedright\arraybackslash}p{0.17\ReportTableWidth}>{\centering\arraybackslash}p{0.16\ReportTableWidth}>{\centering\arraybackslash}p{0.16\ReportTableWidth}>{\centering\arraybackslash}p{0.17\ReportTableWidth}>{\centering\arraybackslash}p{0.17\ReportTableWidth}>{\centering\arraybackslash}p{0.17\ReportTableWidth}@{}}
\caption{人工案例改造前后的 CoreMark 仿真结果}\label{tab:S7}\\
\toprule
{\bfseries Workload} & {\bfseries 改造前周期} & {\bfseries 改造后周期} & {\bfseries 退休指令数} & {\bfseries 改造前 IPC} & {\bfseries 改造后 IPC} \\
\midrule
\endfirsthead
\multicolumn{6}{c}{\small 表~\thetable\quad 人工案例改造前后的 CoreMark 仿真结果（续）}\\[4pt]
\toprule
{\bfseries Workload} & {\bfseries 改造前周期} & {\bfseries 改造后周期} & {\bfseries 退休指令数} & {\bfseries 改造前 IPC} & {\bfseries 改造后 IPC} \\
\midrule
\endhead
\midrule
\multicolumn{6}{r}{\footnotesize 接下页}\\
\endfoot
\bottomrule
\endlastfoot
performance & 2,572,219 & 2,258,952 & 780,059 & 0.303263 & 0.345319 \\[3pt]
validation & 2,747,352 & 2,416,464 & 833,648 & 0.303437 & 0.344987 \\[3pt]
\end{longtable}
\endgroup

这些结果对应条件分支预测、静态直接跳转等整体改造，与第 3 节的两个对照候选分别记录。案例记录至仿真后的工程师评估阶段：工程师要求 Agent 只读评估收益、未覆盖场景及进入物理实现前需要决定的事项。

\subsection{案例呈现与评价范围}

本案例依据开发会话、正式设计、源码和测试结果，展示工程师的架构决定如何转化为可检查的工程约束。文档整理、设计分析和实现分别保留相应 Skill 的调用记录，可沿设计要求追踪到实现位置和验证结果。

\end{document}
